\documentclass[conference]{IEEEtran}

\usepackage[T1]{fontenc}
\usepackage[utf8]{inputenc}
\usepackage[brazil]{babel}
\usepackage{cite}
\usepackage{booktabs}
\usepackage{graphicx}
\usepackage{amsmath}
\usepackage{amssymb}
\usepackage{hyperref}
\usepackage{xurl}
\usepackage{enumitem}

\hypersetup{hidelinks}

\title{Classificacao de Asteroides Potencialmente Perigosos com Analise Explorat\'oria e Machine Learning}
\author{
\IEEEauthorblockN{Alander Menezes Arantes de Avila, Giancarlo Nunes Perli, Gustavo Faria Cardoso \\
Paulo Henrique Perin, Pedro Lucas Ghezzi Bittencourt}
\IEEEauthorblockA{Escola Politecnica, Pontificia Universidade Catolica do Parana (PUCPR) \\
Curitiba, Brasil \\
Autor para correspondencia: menezes.alander@pucpr.edu.br}
}

\begin{document}

\maketitle

\begin{abstract}
Este manuscrito apresenta uma analise de asteroides com foco na classificacao de objetos potencialmente perigosos (PHA). O estudo combina analise exploratoria de dados, avaliacao de variaveis orbitais e fisicas, e modelos de machine learning para predicao binaria de risco. Os resultados indicam que variaveis de geometria orbital sao os principais discriminadores da classe PHA, com desempenho competitivo em cenarios com e sem variaveis definicionais de risco.
\end{abstract}

\begin{IEEEkeywords}
asteroids, potentially hazardous asteroids, binary classification, imbalanced learning, machine learning
\end{IEEEkeywords}

\section{Introducao}
\subsection{Contexto e motivacao}
% TODO: Descrever problema de risco de impacto, contexto astronomico e relevancia da triagem automatizada.

\subsection{Pergunta de pesquisa}
% TODO: Formular pergunta central, por exemplo: quais caracteristicas distinguem PHAs de nao-PHAs?

\subsection{Objetivos}
\begin{itemize}[leftmargin=*]
    \item Objetivo geral: classificar asteroides em PHA e nao-PHA.
    \item Objetivos especificos:
    \begin{itemize}[leftmargin=*]
        \item caracterizar o dataset e suas limitacoes;
        \item identificar variaveis mais informativas;
        \item comparar modelos de classificacao em dados desbalanceados;
        \item discutir trade-offs entre recall e precisao.
    \end{itemize}
\end{itemize}

\section{Dados e Metodologia}
\subsection{Fonte de dados}
% TODO: Inserir fonte Kaggle/NASA, versao, data de acesso e licenca.

\subsection{Descricao do dataset}
\begin{table}[!t]
\centering
\caption{Resumo do dataset utilizado}
\begin{tabular}{ll}
\toprule
Item & Valor \\
\midrule
Total de registros brutos & 958.524 \\
Total de colunas brutas & 45 \\
Registros apos limpeza principal & 938.603 \\
Colunas apos limpeza principal & 38 \\
Classe positiva (PHA=Y) & 2.066 (\~{}0,22\%) \\
Classe negativa (PHA=N) & 936.537 (\~{}99,78\%) \\
\bottomrule
\end{tabular}
\label{tab:dataset_resumo}
\end{table}

\subsection{Limpeza e pre-processamento}
% TODO: Detalhar remocao de colunas identificadoras, tratamento de missing, encoding de variaveis categoricas e split treino/teste.

\subsection{Estrategia para desbalanceamento}
% TODO: Explicar class_weight, metricas adequadas e, se aplicavel, SMOTE.

\subsection{Definicao dos experimentos}
\begin{itemize}[leftmargin=*]
    \item \textbf{Trilha A (com moid\_ld):} experimento com variavel definicional.
    \item \textbf{Trilha B (sem moid\_ld):} experimento sem variavel definicional.
\end{itemize}

\section{Analise Explorat\'oria dos Dados}
\subsection{Analise univariada}
% TODO: Sintetizar os 20 atributos analisados e destacar os mais relevantes.

\subsection{Analise multivariada}
% TODO: Descrever hipoteses e achados principais (e.g., H x MOID, a x e, correlacoes de Spearman).

\subsection{Principais achados da EDA}
\begin{itemize}[leftmargin=*]
    \item Geometria orbital discrimina melhor PHAs.
    \item Forte multicolinearidade entre variaveis orbitais derivadas.
    \item Diametro e albedo apresentam alta ausencia e menor poder discriminativo.
\end{itemize}

\begin{figure}[!t]
    \centering
    % \includegraphics[width=\columnwidth]{figuras/final_plot_1.png}
    \fbox{\parbox[c][4.4cm][c]{0.95\columnwidth}{\centering Placeholder: Figura 1 (Regiao de risco PHA)}}
    \caption{Regiao conjunta dos criterios de risco para classificacao PHA.}
    \label{fig:final_plot_1}
\end{figure}

\begin{figure}[!t]
    \centering
    % \includegraphics[width=\columnwidth]{figuras/final_plot_2.png}
    \fbox{\parbox[c][4.4cm][c]{0.95\columnwidth}{\centering Placeholder: Figura 2 (Orbitas por classe)}}
    \caption{Comparacao qualitativa de orbitas para PHA e nao-PHA.}
    \label{fig:final_plot_2}
\end{figure}

\section{Modelagem de Machine Learning}
\subsection{Modelos avaliados}
\begin{itemize}[leftmargin=*]
    \item Regressao Logistica
    \item Random Forest
    \item XGBoost
\end{itemize}

\subsection{Protocolo de avaliacao}
% TODO: Explicar hold-out estratificado, validacao cruzada (se usada) e metricas principais.

\subsection{Resultados quantitativos}
\begin{table*}[!t]
\centering
\caption{Comparacao de desempenho por trilha e modelo}
\footnotesize
\begin{tabular}{llllll}
\toprule
Trilha & Modelo & Precision (PHA) & Recall (PHA) & F1 (PHA) & PR-AUC \\
\midrule
A & Regressao Logistica & 0,24 & 1,00 & 0,39 & 0,6678 \\
A & Random Forest & 0,98 & 0,99 & 0,98 & 0,9974 \\
A & XGBoost & 0,92 & 0,92 & 0,92 & 0,9748 \\
B & Regressao Logistica & 0,16 & 1,00 & 0,27 & 0,3090 \\
B & Random Forest & 0,75 & 0,39 & 0,52 & 0,6737 \\
B & XGBoost & 0,57 & 0,82 & 0,67 & 0,7642 \\
\bottomrule
\end{tabular}
\label{tab:ml_resultados}
\end{table*}

\subsection{Interpretacao dos resultados}
% TODO: Discutir leakage semantico na Trilha A, valor cientifico da Trilha B e trade-off recall x precisao.

\section{Discussao}
\subsection{Implicacoes praticas}
% TODO: Relacionar achados com priorizacao de observacao e monitoramento de risco.

\subsection{Limitacoes}
\begin{itemize}[leftmargin=*]
    \item Alto percentual de missing em variaveis fisicas (diameter/albedo).
    \item Dependencia potencial de variaveis definicionais em cenarios com moid\_ld.
    \item Possiveis mudancas de distribuicao temporal do catalogo.
\end{itemize}

\subsection{Ameacas a validade}
% TODO: Inserir validade interna, externa e de construcao.

\section{Conclusao}
% TODO: Consolidar 3--5 mensagens finais e responder explicitamente a pergunta de pesquisa.

\section*{Agradecimentos}
% TODO: Agradecimentos ao professor, instituicao e eventuais apoios.

\begin{thebibliography}{99}

\bibitem{nasa_jpl}
NASA/JPL, "Small-Body Database and NEO Resources," [Online]. Available: \url{https://ssd.jpl.nasa.gov/}. Accessed: Jun. 4, 2026.

\bibitem{kaggle_asteroids}
Sakhawat, "Asteroid Dataset," Kaggle, [Online]. Available: \url{https://www.kaggle.com/datasets/sakhawat18/asteroid-dataset/data}. Accessed: Jun. 4, 2026.

\bibitem{imbalanced_learning}
H. He and E. A. Garcia, "Learning from imbalanced data," \textit{IEEE Trans. Knowl. Data Eng.}, vol. 21, no. 9, pp. 1263--1284, Sep. 2009.

\end{thebibliography}

\appendix
\section{Checklist de reproducibilidade}
Checklist completo salvo no arquivo auxiliar \textit{CHECKLIST-MANUSCRITO-IEEE.md}.

Itens minimos para reproducibilidade:
\begin{itemize}[leftmargin=*]
    \item Versoes de bibliotecas e ambiente.
    \item Seeds aleatorias utilizadas.
    \item Ordem de execucao dos experimentos.
    \item Vinculo direto entre resultados no texto e outputs dos notebooks finais.
\end{itemize}

\section{Material suplementar}
Guia de revisao e submissao salvo em \textit{GUIA-EXECUCAO-MANUSCRITO.md}.

Conteudos recomendados para o material suplementar:
\begin{itemize}[leftmargin=*]
    \item Tabelas adicionais por modelo e por trilha experimental.
    \item Matrizes de confusao completas no conjunto de teste.
    \item Curvas PR/ROC e analise de limiar de decisao.
    \item Configuracoes de hiperparametros e tempo de execucao.
\end{itemize}

\end{document}
